\begin{abstract}
An orchestrating agent often delegates work to cheaper sub-agents and must choose a model and a reasoning effort level for each, which together drive cost, wall-clock time, and answer quality. We measure how much these choices matter for Claude Haiku~5.5 and Claude Sonnet~5.5 run as Claude Code sub-agents in terminal panes. Three experiments covered 229 tool-free reasoning runs, 120 coding runs on an easy set of 12 tasks, and 150 coding runs on a hard set of 10 tasks, all built from a private Python repository, with API-equivalent costs computed from session transcripts. Haiku \texttt{low} passed all 30 fix and feature runs across both sets (Wilson 95\,\% interval 89--100\,\%), at \$0.041 per success over all hard-set tasks, reviews included, against \$0.207 for Sonnet \texttt{low}. Effort above \texttt{medium} added cost with no accuracy gain we could detect, though nearly every task passed at every configuration: Sonnet \texttt{xhigh} cost 2.2 to 2.6 times Sonnet \texttt{low} per success. Model and effort did matter for long exact reasoning without tools: Sonnet \texttt{xhigh} alone traced 1,300 statements without failure, while Haiku \texttt{high} and \texttt{xhigh} took a wrong shortcut at 40. We hypothesize that code-review misses are shared within a model more than across models; this rests on a handful of cases, with exceptions in both directions. We route lookups and code changes to Haiku \texttt{low}, review to both models in parallel, and long tool-free reasoning to Sonnet \texttt{xhigh}. The main caveat is a ceiling effect: nearly every coding task passed at every configuration, so we did not find where coding reliability breaks down.
\end{abstract}
